% Epilogue of The Hydrogen Atoms of Deep Learning.
\specialchapter{Epilogue}{Fine structure: what the solvable models get wrong together}{100}\label{ch:epilogue}
\begin{refsection}[chapters/17-epilogue.bib]

Sixteen models, each solved in its own chapter, each with its own list of failures. Read one at a time, they look like sixteen separate tricks. Read side by side, they share a small number of mechanisms, and they also share their failures. A failure that one model has is a fact about that model. A failure that every model has is a fact about the method, and it is the fine structure that the next generation of models has to explain. This epilogue collects both, lists phenomena that still have no hydrogen atom, and ends with a recipe for building one.

\addsec{Mechanisms that recur}

The table below lists the mechanisms that carry more than one chapter, the chapters in which they do the main work, and what each one computes. None of them is new to this book; most are older than deep learning. What is striking is how few there are.

\begin{fullwidth}
\small
\renewcommand{\arraystretch}{1.22}
\begin{tabularx}{\linewidth}{@{}>{\raggedright}p{0.2\linewidth}>{\raggedright}p{0.14\linewidth}X@{}}
\toprule
\textsc{mechanism} & \textsc{chapters} & \textsc{what it computes}\\
\midrule
Escape from a saddle & 2, 3, 7, 8, 9, 10, 13, 14 & The length of a plateau: the logarithm of the initial distance from the saddle, divided by the unstable rate\\
Conservation laws from rescaling symmetries & 2, 3, 13, 15 & Which minimizer training reaches, how sharp it is, and which initializations make the dynamics closed\\
Reduction to order parameters & 1, 5, 7, 8, 13 & A closed recursion or ordinary differential equation in a few overlaps, from rotational invariance of Gaussian data\\
Harmonic analysis & 5, 6, 7, 9, 10 & What the gradient sees at each degree: Hermite coefficients, spherical harmonics, characters of $\Z_p$ and of the cube\\
Margins as attractors & 4, 10, 12 & The direction that an exponentially tailed loss selects after the training error reaches zero\\
Spectra of random and structured matrices & 1, 6, 13, 16 & Risk as a functional of a spectrum: Marchenko--Pastur, kernel eigenvalues, power-law tails\\
A limit that closes the equations & 1, 5, 6, 7, 8, 16 & Width, dimension or sample size taken to infinity in a fixed ratio or a fixed order\\
\bottomrule
\end{tabularx}
\end{fullwidth}

The first row deserves a comment, because it is the most reliable prediction in the book. Near a saddle the dynamics is linear, and an unstable mode with rate $\lambda$ grows like $e^{\lambda t}$. Starting at distance $\eps$, it needs a time $\lambda^{-1}\log(1/\eps)$ to reach order one. This is the normal-mode theorem of \hyperref[ch:prologue]{the prologue} with the sign of one eigenvalue reversed. It gives the half-learning time $\log(s/a_0)/(2s)$ of a two-layer linear network (\cref{ch:deeplinear}), the time $\log(1/\alpha)/(2\abs{w_{\star,i}})$ at which a diagonal network learns coordinate $i$ (\cref{ch:diagonal}), the $\log d$ plateau of the committee machine (\cref{ch:committee}), the delay $\lambda^{-1}\log(\alpha/\alpha_c)$ of grokking under weight decay (\cref{ch:grokking}), and the escape times of linear attention from small initialization (\cref{ch:icl,ch:induction}). The single-index model and the sparse parity (\cref{ch:singleindex,ch:parity}) are the same story at a degenerate saddle, where the unstable direction grows like a power rather than an exponential, and the logarithm becomes a power of the dimension. Whenever training shows a plateau followed by a drop, this is the first model to try, and the two numbers to measure are the initial distance and the unstable rate.

The second row is less visible but as useful. A rescaling symmetry of the parameterization, such as $(W_{\ell+1},W_\ell)\mapsto(W_{\ell+1}A^{-1},AW_\ell)$ for a linear network or $(u,v)\mapsto(cu,v/c)$ for a product, gives a conserved quantity under gradient flow: the balancedness of \cref{ch:deeplinear}, the products $u_iv_i=\alpha^2$ of \cref{ch:diagonal}. Under gradient descent with a finite step the same quantity is no longer conserved but evolves by an exact law, and in \cref{ch:stability} that law is what forces the catapult. In every case the conserved or slowly varying quantity selects among minimizers that the loss cannot tell apart.

The third and fourth rows are the reason the models are solvable at all. Gaussian data are rotationally invariant, so a high-dimensional state collapses to a few overlaps; uniform data on a group are invariant under translation, so the right basis is the group's characters. The fifth row is the mechanism of implicit bias in classification: once the data are separated, an exponentially tailed loss is dominated by its smallest margins, and the trajectory follows the maximum-margin problem. Three chapters of very different appearance, on logistic regression, modular addition and the last layer of a classifier, are three solutions of that problem.

\addsec{What the models get wrong together}

Each chapter ends with what its model gets wrong. Putting the lists side by side, six failures recur.

\paragraph{The data are isotropic, independent, or uniform.} Most of the sixteen models assume Gaussian inputs with identity covariance, uniform inputs on a group or a cube, or independent features. This is what makes the order parameters close. Real data have structure at every scale, and in the models that structure enters, if at all, only through a spectrum (\cref{ch:doubledescent,ch:scaling}) or through a target that is given rather than learned. Where a chapter tests what happens without isotropy, the prediction changes: anisotropic inputs change the information exponent (\cref{ch:singleindex}), and structured inputs change or remove the plateau of the committee machine (\cref{ch:committee}).

\paragraph{One mechanism at a time.} Every model isolates one phenomenon and removes everything that could compete with it. In a trained network the mechanisms run together. The delay of grokking in \cref{ch:grokking} had to be modelled by a different network from the one that explains the Fourier circuit; the induction head of \cref{ch:induction} is one circuit among many that form during the same phase change; implicit bias, edge-of-stability dynamics and feature learning act on the same weights at the same time. No model in this book solves an interaction between two of its own mechanisms.

\paragraph{Limits in a convenient order.} Solvability comes from a limit, and the limit is usually taken in the order that makes the equations close: width to infinity at fixed depth (\cref{ch:signalprop,ch:lazy}), dimension to infinity before time (\cref{ch:singleindex,ch:committee}), time to infinity for a margin (\cref{ch:margin}). The interesting behavior often lives where two limits compete. Signal propagation fails when depth is comparable to width; the committee machine's plateau length is set by fluctuations of order $d^{-1/2}$ that its own ordinary differential equation omits; the margin converges at a rate in $1/\log t$ that no training run reaches.

\paragraph{Gradient flow, or plain gradient descent.} Almost every theorem is proved for gradient flow, online stochastic gradient descent with vanishing steps, or full-batch gradient descent with a small step. Only \cref{ch:stability} takes a large step seriously, and no chapter treats an adaptive optimizer. Yet the chapters themselves record that the optimizer changes the answer: step size and label noise change the implicit bias of diagonal networks (\cref{ch:diagonal}), steepest descent and Adam change the margin that classification selects (\cref{ch:margin}), and reusing samples beats the information exponent (\cref{ch:singleindex}).

\paragraph{Special initializations.} Several solutions require an initialization chosen to make the dynamics solvable: aligned and balanced for deep linear networks, symmetric for the committee machine and the sparse parity, small for diagonal networks and attention. The chapters show that random initializations keep the scaling and change the constants, but the exact solutions are exact only on a set of measure zero.

\paragraph{Minimal architectures.} The models are linear, diagonal, two-layer, or one layer of attention. Residual connections, normalization layers and softmax attention appear only in the sections on what the models get wrong. The transformer simulation of \cref{ch:induction} did not reach the bigram level with a standard linear readout and needed a modified one.

\medskip

The book's own simulations are a fair test of these failures, because every chapter compares its model with a finite simulation, and several of them disagreed with the model in an instructive way. A small ReLU network trained in \cref{ch:collapse} collapsed only partially. The networks of \cref{ch:grokking} are far from the maximum-margin solutions that explain their Fourier features. The fitted exponents of language models violate the relation between data and parameter exponents that both models of \cref{ch:scaling} predict. And the transformers that match least squares on regression prompts cannot be performing the single gradient step that \cref{ch:icl} proves optimal for one layer.

The pattern across these comparisons is consistent. What the models predict well are scalings and orders of events: a time that grows like $\log(1/\alpha)$ or a power of $d$, a threshold at $\eta\Lambda=2$, the order in which modes are learned, the geometry of an optimum. What they predict badly are constants, and which mechanism is active in a given network. This is the same division as in physics, where the hydrogen atom gets the structure of the spectrum right and the fine structure wrong, and where the corrections are small. In deep learning the corrections are not always small.

\addsec{Missing hydrogen atoms}

Several phenomena that matter in practice have no model in this book, either because none exists or because the existing ones do not yet pass the four tests of \hyperref[ch:prologue]{the prologue}.

\emph{Adaptive optimizers.} Almost every large network is trained with Adam \parencite{KingmaBa2015} or a relative of it. Sign descent and steepest descent in other norms have partial theories (\cref{ch:margin}), but there is no solvable model that predicts when Adam trains a network that gradient descent cannot, or the step-size schedules used in practice.

\emph{Progressive sharpening.} The model of \cref{ch:stability} explains why gradient descent at a large step does not diverge, but not why the sharpness rises to $2/\eta$ in the first place, which \textcite{CohenEtAl2021} observed in many networks. A model that produces sustained sharpening from its own dynamics is missing.

\emph{Normalization layers.} Batch normalization \parencite{IoffeSzegedy2015} and layer normalization \parencite{BaKirosHinton2016} change the geometry of the loss, make it invariant to the scale of the weights, and interact with weight decay and the step size. They appear in this book only as things the models leave out.

\emph{Data mixtures and curricula.} Real training sets are mixtures of sources with different difficulty and redundancy, and the order in which examples arrive can matter \parencite{BengioEtAl2009}. The models here draw every example from one fixed distribution.

\emph{Learning from feedback.} Reinforcement learning from human preferences \parencite{ChristianoEtAl2017,OuyangEtAl2022} trains a pretrained model on a reward that is itself learned. Its dynamics, including how far the model moves from its pretrained state, have no solvable model of the kind collected here.

\emph{Chain of thought.} Prompting a model to write intermediate steps improves its answers on reasoning tasks \parencite{WeiEtAl2022}. The sparse parity of \cref{ch:parity} suggests why intermediate steps can help a learner, since they break a high-degree target into low-degree pieces, but a model that predicts the size of the effect in a trained transformer does not exist yet.

\emph{Long-context retrieval.} Language models use information at the start and end of a long input better than information in the middle \parencite{LiuEtAl2024}. The models of \cref{ch:icl,ch:induction} treat every position in the prompt alike.

\addsec{A recipe for a new hydrogen atom}

The chapters suggest a procedure, which a reader can apply to each of the phenomena above.

\begin{enumerate}
\item \emph{Choose a phenomenon with a number.} An exponent, a time scale, a threshold or a geometry, observed robustly in several networks. Without a number the model cannot fail, and a model that cannot fail cannot teach.
\item \emph{Find the symmetry.} Look for the invariance that reduces the state to a few variables: rotations of Gaussian inputs, translations on a group, a rescaling of the parameters. If no symmetry is visible, build it into the data.
\item \emph{Choose the limit and the initialization.} Take the limit (dimension, width, time, initialization scale) and choose the starting point that makes the reduced equations close. Write down what this choice throws away; it will be the first section of ``What the model gets wrong''.
\item \emph{Solve the smallest instance completely.} Find the saddle, the conserved quantity or the margin problem. Prove the main theorem, and label as heuristic whatever you cannot prove.
\item \emph{Predict a number and simulate with no fitted constant.} Compare the prediction with a finite simulation of the model, and then with a network that the model does not describe. The disagreement with the second is the result.
\item \emph{Break one assumption on purpose.} Change the data, the optimizer, the architecture or the limit, one at a time, and solve again. This is what the variation exercises of every chapter ask for, and it is how one model becomes the next.
\item \emph{Write the model card.} If the six lines cannot be filled honestly, in particular the line on what the model gets wrong, the model is not ready.
\end{enumerate}

The hydrogen atom became the foundation of atomic physics not because it was right, but because every way in which it was wrong pointed to something real: relativity, the electron's spin, the quantized field. The models in this book are wrong in ways that are known and specific. The hope is that their failures, studied as carefully as their solutions, point to what deep learning theory needs next.

\notesandrefs

The table of mechanisms condenses material from the chapters, where each mechanism is proved and referenced. The observation that the time to escape a saddle scales like $\log(1/\eps)/\lambda$ is classical in the theory of dynamical systems, and its role in the chapters is discussed in the notes to \cref{ch:deeplinear,ch:diagonal,ch:committee}. Adam is due to \textcite{KingmaBa2015}; batch and layer normalization to \textcite{IoffeSzegedy2015} and \textcite{BaKirosHinton2016}; curriculum learning was named by \textcite{BengioEtAl2009}; learning from human preferences is due to \textcite{ChristianoEtAl2017} and was applied to language models by \textcite{OuyangEtAl2022}; chain-of-thought prompting is due to \textcite{WeiEtAl2022}; the position dependence of long-context retrieval was measured by \textcite{LiuEtAl2024}; progressive sharpening was named by \textcite{CohenEtAl2021}.

\printchapterbib
\end{refsection}
